%! AP Statistics — Unit 1, Lesson 1.8 — Group Activity ANSWER KEY
\documentclass[10pt]{article}
\usepackage{apstats-article}
\usepackage{apstats-key}

%=============================================================================
\begin{document}

\pageheader{Unit 1, Lesson 1.8}{Group Activity --- Answer Key}
\begin{tabularx}{\linewidth}{X X X}
Name:\;\ans{Key} & Date:\; & Period:\;
\end{tabularx}

\vspace{0.10in}

\begin{tcolorbox}[colback=sky, colframe=navy, boxrule=0.8pt, arc=2mm,
    left=3mm, right=3mm, top=2mm, bottom=2mm]
\small
\textbf{Part 1} (group, 8 min): Compute the five-number summary and check for
outliers.\\
\textbf{Part 2} (group, 7 min): Interpret a pre-drawn modified boxplot.\\
\textbf{Part 3} (individual, 5 min): Write a SOCS description and compare displays.
\end{tcolorbox}

\vspace{0.10in}

{\large\bfseries\color{navy} Part 1 \quad Compute the Five-Number Summary}

\vspace{0.08in}

\noindent\small Data (sorted, $n=17$):\quad
52, 56, 58, 60, 62, 63, 64, 66, \textbf{68}, 70, 72, 73, 75, 78, 80, 84, 110

\vspace{0.10in}

\begin{enumerate}[leftmargin=1.5em, itemsep=6pt]

  \item $n = $ \ans{17} values. Median is the \ans{9th} value. \quad $M = $ \ans{68} bpm

  \item Lower half (first 8): \ans{52, 56, 58, 60, 62, 63, 64, 66}\\[4pt]
        $Q_1 = (62+63)/2 = $ \ans{62.5} bpm \hfill
        Upper half (last 8): \ans{70, 72, 73, 75, 78, 80, 84, 110}\\[4pt]
        $Q_3 = (75+78)/2 = $ \ans{76.5} bpm \hfill
        $\text{IQR} = 76.5 - 62.5 = $ \ans{14} bpm

  \item Five-number summary:\\[4pt]
        \renewcommand{\arraystretch}{1.5}
        \begin{tabularx}{0.85\linewidth}{@{} Y Y Y Y Y @{}}
        \rowcolor{sky}
        $\min$ & $Q_1$ & $M$ & $Q_3$ & $\max$ \\
        \midrule
        \ans{52} & \ans{62.5} & \ans{68} & \ans{76.5} & \ans{110} \\
        \end{tabularx}

  \item Lower fence $= 62.5 - 1.5(14) = 62.5 - 21 = $ \ans{41.5 bpm}\\[4pt]
        Upper fence $= 76.5 + 1.5(14) = 76.5 + 21 = $ \ans{97.5 bpm}\\[4pt]
        Is 110 an outlier? \ans{Yes} \quad Because: \ans{$110 > 97.5$ (upper fence)}

  \item Last non-outlier $= $ \ans{84} bpm.

\end{enumerate}

\vspace{0.10in}

{\large\bfseries\color{navy} Part 2 \quad Interpret the Modified Boxplot}

\vspace{0.08in}

\begin{center}
\begin{tikzpicture}[scale=1.1]
  \draw[thick] (0.4,0) -- (8.0,0);
  \foreach \v/\lbl in {0.50/50, 1.50/60, 2.50/70, 3.50/80,
                        4.50/90, 5.50/100, 6.50/110, 7.50/120} {
    \draw (\v,-0.12) -- (\v,0.12);
    \node[below, font=\small] at (\v,-0.12) {\lbl};
  }
  \node[below, font=\small\itshape] at (4.0,-0.54) {Pulse rate (bpm)};
  \draw[thick, navy] (0.70, 0.35) -- (1.75, 0.35);
  \draw[thick, navy] (1.75, 0.10) rectangle (3.15, 0.60);
  \draw[thick, navy] (2.30, 0.10) -- (2.30, 0.60);
  \draw[thick, navy] (3.15, 0.35) -- (3.90, 0.35);
  \draw[thick, navy] (0.70, 0.20) -- (0.70, 0.50);
  \draw[thick, navy] (3.90, 0.20) -- (3.90, 0.50);
  \filldraw[navy] (6.50, 0.35) circle (0.10);
\end{tikzpicture}
\end{center}

\begin{enumerate}[leftmargin=1.5em, itemsep=6pt, resume]

  \item $Q_1 = $ \ans{62.5}, $M = $ \ans{68}, $Q_3 = $ \ans{76.5}. \\[3pt]
        \ans{Yes, these match the calculated values.}

  \item Value of outlier dot: \ans{110 bpm}. \quad
        Right whisker end represents: \ans{the last non-outlier value (84 bpm)}.

  \item Between 62.5 and 76.5 bpm: $\approx$ \ans{50\%}. \quad
        Between 52 and 68 bpm: $\approx$ \ans{50\%}.

\end{enumerate}

\vspace{0.10in}

{\large\bfseries\color{navy} Part 3 \quad SOCS Description \& Comparing Displays}

\begin{enumerate}[leftmargin=1.5em, itemsep=6pt, resume]

  \item \textbf{Shape:} \ans{The distribution of pulse rates is slightly right-skewed.
        The right whisker (84 − 76.5 = 7.5 bpm) is longer than the left
        (62.5 − 52 = 10.5 bpm), and there is a high outlier at 110 bpm that
        extends the upper tail.}\\[5pt]
        \textbf{Outlier(s):} \ans{There is one outlier at 110 bpm. This value exceeds
        the upper fence of 97.5 bpm ($= 76.5 + 1.5 \times 14$) and is plotted
        separately as a dot on the modified boxplot.}\\[5pt]
        \textbf{Center:} \ans{The median pulse rate for these student athletes is
        68 bpm, meaning half of the athletes had pulse rates below 68 bpm and
        half had pulse rates above.}\\[5pt]
        \textbf{Spread:} \ans{The IQR is 14 bpm ($Q_3 - Q_1 = 76.5 - 62.5$),
        meaning the middle 50\% of pulse rates fall within a 14-bpm window.
        Excluding the outlier, the range is $84 - 52 = 32$ bpm.}

  \item \ans{The modified boxplot clearly shows the presence of one outlier
        (110 bpm) as a separate dot, distinguishing it from the bulk of the data.}

  \item \ans{A histogram would reveal the shape of the distribution in more
        detail --- for example, whether the data cluster near 65--70 bpm or
        whether there is a gap between 84 and 110 bpm.}

\end{enumerate}

\begin{teachernote}
\textbf{Grading notes:}
\begin{itemize}[itemsep=2pt]
  \item Part 1 (Q1--Q5): 1 pt each; require IQR=14, fences correct, 110 identified as outlier.
  \item Part 2 (Q6--Q8): Q6 checks reading; Q7 requires naming "last non-outlier" not "fence"; Q8: 50\% for both (students may say "approximately 25\%" for each quarter — accept if reasoning is shown).
  \item Part 3 (Q9): SOCS: 1 pt per component. Shape must reference whisker asymmetry. Outlier must give fence calculation. Center must state median and interpret. Spread must give IQR in context.
  \item Q10--Q11: Accept any valid answer with explanation.
\end{itemize}
\end{teachernote}

\end{document}
